%
% 34-singulaer.tex -- Singularitäten
%
% (c) 2026 Prof Dr Andreas Müller
%
\subsection{Singuläre Punkte}
Die Laurent-Reihenentwicklung ermöglicht, eine Funktion in der Umgebung
eines singulären Punktes genauer zu untersuchen.
Sei $U$ ein offene Teilmenge von $\mathbb{C}$ und $z_0$ ein isolierter
Punkt von $\mathbb{C}\setminus U$.
Ein auf $U$ definierte, holomorphe Funktion $f\colon U\to\mathbb{C}$
ist also in jeder genügend kleinen Umgebung von $z_0$ in jedem Punkt
definiert ausser im Punkt $z_0$.
Eine solche Stelle heisst ein \emph{isolierter} singulärer Punkt.
\index{isoliert}%
\index{singularer Punkt@singulärer Punkt}%

Die Funktion $f$ ist in jedem Ringgebiet $R(z_0;r_1,r_2)$, um den Punkt
$z_0$ definiert, sofern nur $r_1$ klein genug ist und für 
alle inneren Radien $r_1<r_2$.
Nach Satz~\ref{buch:analytisch:laurent:satz:laurentreihe} gibt es
eine Laurent-Reihenentwicklung
\index{Laurent-Reihe}%
\begin{equation}
f(z)
=
\sum_{k=0}^\infty a_k(z-z_0)^k
+
\sum_{k=1}^\infty b_k(z-z_0)^{-k},
\label{buch:analytisch:laurent:eqn:laurentsing}
\end{equation}
die für $z$ beliebig nahe an $z_0$ konvergiert.
Es folgt, dass die Reihe
\[
u(z)
=
\sum_{k=1}^\infty b_kz^k
\]
für beliebig grosse $z$ konvergiert, sie ist also in ganz $\mathbb{C}$
definiert.
Am Nullpunkt ist $u(0)=0$.
Die zweite Summe in der Laurent-Reihe
\eqref{buch:analytisch:laurent:eqn:laurentsing}
kann daher als $u(1/(z-z_0))$ geschrieben werden.
Die Funktion $u(1/(z-z_0))$ heisst der \emph{singuläre Teil} der Funktion
$f$ in einer Umgebung von $z_0$.
\index{singularer Teil@singulärer Teil}%

Falls $u=0$ ist, dann handelt es sich um eine \emph{hebbare} Singularität,
\index{hebbar}%
\index{Singularitat@Singularität}%
\index{Singularitat, hebbar@Singularität, hebbar}%
die Funktion kann stetig in den Punkt $z_0$ hinein fortgesetzt werden,
die zweite Summe in \eqref{buch:analytisch:laurent:eqn:laurentsing}
ist nicht nötig.

\begin{definition}[Ordnung eines singlären Punktes]
Sei $z_0$ ein singulärer Punkt der Funktion $f(z)$ und $u(z)$ der 
singuläre Teil an der Stelle $z_0$.
Wenn $u(z)$ ein Polynom vom Grad $n$ ist, dann heisst $z_0$ ein
\emph{Pol der Ordnung $n$}.
\index{Pol}%
\index{Ordnung!Pol}%
Falls $u(z)$ kein Polynom ist, heisst $z_0$ eine \emph{essentielle
Singularität}.
\index{essentielle Singularitat@essentielle Singularität}%
Falls $u=0$ ist und $a_k=0$ für $k\le n$ in 
\eqref{buch:analytisch:laurent:eqn:laurentsing},
dann heisst $z_0$ eine \emph{Nullstelle der Ordnung $n$}.
\index{Ordnung!Nullstelle}%
Die Funktion $\omega(z_0,f)$ ist definiert durch
\[
\omega(z_0,f)
=
\begin{cases}
-\infty&\qquad\text{falls $z_0$ eine essentielle Singularität von $f$ ist}\\
-n&\qquad\text{falls $z_0$ ein Pol der Ordnung $n$ von $f$ ist}\\
\phantom{-}m&\qquad\text{falls $z_0$ eine Nullstelle der Ordnung $m$ von $f$ ist}\\
\phantom{-}\infty&\qquad\text{falls $f=0$ ist}
\end{cases}
\]
\end{definition}

Für die Ordnung $\omega(z_0,f)$ gelten die folgenden Rechenregeln.
\begin{align*}
\omega(z_0,f+g) &\ge \min\bigl(\omega(z_0,f),  \omega(z_0,g)\bigr)
\\
\omega(z_0,fg) &= \omega(z_0,f) + \omega(z_0,g).
\end{align*}

Wenn $f$ eine Singularität der Ordnung $n$ an der Stelle $z_0$ hat, 
dann ist $(z-z_0)^{-n}f(z)$ eine analytische Funktion mit einer
hebbaren Singularität an der Stelle $z_0$, die dort nicht verschwindet.
Insbesondere ist dann auch $1/(z-z_0)^{-n}\cdot 1/f(z)$ eine in einer
Umgebung von $z_0$ holomorphe Funktion, so dass man
\[
\omega(z_0,f)
=
-\omega(z_0,1/f)
\]
schliessen kann.
